% auto-generated by scripts/altloading.mjs
\begin{table}[htbp]
  \centering
  \caption{Alternative eager-loading sensitivity on the deep fetch:
    policy-selected documented throughput (req/s) versus the alternative loading strategy, both
    engines, for the layers whose alternative is a byte-identical drop-in. The table reports only byte-identical drop-ins; excluded alternatives remain
    configuration-portability findings rather than timed treatments. This is a secondary
    run that predates the accepted primary campaign, so its selected column sits up to
    about $14\%$ below the corresponding primary median; only the within-table
    selected-versus-alternative contrast is interpreted, never the absolute
    values. On both engines the policy-selected documented path is the \emph{faster} strategy for \texttt{sequelize} and \texttt{mikroorm}, so their deep-fetch deficit does not appear to be an artifact of a poor selected strategy. \texttt{typeorm} is omitted: its alternative strategy was not a byte-identical drop-in in the pinned versions (TypeORM's query strategy errors on the ordered deep-fetch \texttt{findOne}; Objection's \texttt{withGraphJoined} changes row handling), which is itself a portability caveat.}
  \label{tab:altloading}
  \begin{tabular}{l l r r r r}
    \toprule
    & & \multicolumn{2}{c}{PostgreSQL} & \multicolumn{2}{c}{MySQL} \\
    \cmidrule(lr){3-4}\cmidrule(lr){5-6}
    Layer & Switch & selected & alt. & selected & alt. \\
    \midrule
    \texttt{sequelize} & join$\to$select-in & 915 & 690 & 830 & 602 \\
    \texttt{mikroorm} & join$\to$select-in & 447 & 357 & 420 & 324 \\
    \texttt{objection} & select-in$\to$join & -- & -- & 820 & 1301 \\
    \bottomrule
  \end{tabular}
\end{table}
